% Corrected French for printed p.169; departures are append-only below.
\SourcePage{169}

\LayerTarget{tohoku:proof:R-quotient-ZYm:p2}
\noindent nulle. D'autre part, la restriction de $R_m/R_{m-1}$ à
$Y_m=U_m\cap \mathrm C V_{m-1}$ est isomorphe à celle de $R_m$ (puisque
celle de $R_{m-1}$ est nulle) et y est engendrée par la restriction de la
section $f_m$, donc elle est isomorphe au faisceau constant $Z$.

\LayerTarget{tohoku:proof:3.6.1:p3}
Ceci prouve bien que $R_m/R_{m-1}$ est isomorphe à $Z_{Y_m}$, et nous ramène
à prouver que $T^{p+1}(Z_Y)\in C''$ si $Y$ est une partie localement fermée
de $X$. On a alors $Y=V\cap \mathrm C U$, où $U$ et $V$ sont deux parties
ouvertes de $X$. Soit $W$ la partie ouverte de $X$ réunion de $U$ et $V$, on
a une suite exacte évidente :
\[
0\longrightarrow Z_U\longrightarrow Z_W\longrightarrow Z_Y
\longrightarrow0,
\]
d'où une suite exacte
\[
T^{p+1}(Z_W)\longrightarrow T^{p+1}(Z_Y)
\longrightarrow T^{p+2}(Z_U).
\]
Par hypothèse, les termes extrêmes de cette suite exacte sont dans $C''$, il
en est donc de même de $T^{p+1}(Z_Y)$, ce qui achève la démonstration de la
proposition.

\LayerTarget{tohoku:proposition:3.6.2}
\noindent\textsc{Proposition 3.6.2.} --- \textit{Soient $C$, $C'$ deux
catégories abéliennes satisfaisant l'axiome AB 5) (cf. 1.5), on suppose que
$C$ admet un générateur (cf. 1.9). Soit $T$ un foncteur covariant de $C$ dans
$C'$. Pour que les foncteurs dérivés droits $R^pT$ commutent aux limites
inductives, il suffit que :}
\begin{description}
\LayerTarget{tohoku:subitem:3.6.2:a}
\item[\textit{a)}] \textit{$T$ permute aux limites inductives ;}
\LayerTarget{tohoku:subitem:3.6.2:b}
\item[\textit{b)}] \textit{Si $M=\varinjlim_i M_i$ dans $C$, où les $M_i$
sont injectifs, alors $M$ est $T$-acyclique, i.\,e.
$R^pT(M)=0$ pour $p>0$.}
\end{description}

\begin{quote}
\small
\textbf{Qualification éditoriale.} Avec la définition de la p.~143, pour un
foncteur covariant qui n'est supposé qu'additif, le morphisme canonique
$T(M)\longrightarrow R^0T(M)$ n'est pas nécessairement un isomorphisme.
L'annulation $R^qT(M)=0$ pour $q>0$ exigée en b) ne fournit donc pas, à
elle seule, la comparaison en degré zéro requise pour l'égalité
$R^pT(A)=H^p(T(C))$ employée plus bas. La proposition et sa démonstration
sont valables si l'on ajoute l'hypothèse que $T$ est exact à gauche.
Sans cette hypothèse, il faut renforcer b) en demandant, pour tout
$M=\varinjlim_i M_i$ qui y figure, à la fois que
$T(M)\longrightarrow R^0T(M)$ soit un isomorphisme et que
$R^qT(M)=0$ pour $q>0$. Cette qualification ne modifie ni a), ni la
construction des résolutions compatibles, ni les isomorphismes de
comparaison obtenus sous l'une ou l'autre des hypothèses ainsi précisées.
\end{quote}

\LayerTarget{tohoku:proof:3.6.2:p1}
(La condition b) est évidemment nécessaire pour que la conclusion soit
valable, et de même a) si $T$ est exact à gauche donc $R^0T=T$). Notons
d'abord que les $R^pT$ sont bien définis, en vertu de th. 1.10.1. Soit
$(A_i)_{i\in I}$ un système inductif dans $C$, $A$ sa limite inductive, on veut
prouver que les morphismes naturels

\LayerTarget{tohoku:eq:derived-colimit-comparison}
\[
\varinjlim_i R^pT(A_i)\longrightarrow R^pT(A)
\]
sont bijectifs.

\LayerTarget{tohoku:construction:inductive-injective-resolutions}
On va montrer d'abord qu'il existe un système inductif $(C_i)_{i\in I}$ de
complexes (à valeurs dans $C$) et une ``augmentation''
$(A_i)\longrightarrow(C_i)$, de telle façon que pour tout $i\in I$,
$A_i\longrightarrow C_i$ soit une résolution droite de $A_i$ par des objets
injectifs.

\LayerTarget{tohoku:construction:diagram-category-resolution}
En effet, considérons la catégorie $I(C)$ des systèmes inductifs sur $I$ à
valeurs dans $C$, c'est une catégorie de diagrammes, qui d'après prop. 1.6.1
satisfait aux mêmes hypothèses que $C$. En vertu de théorème 1.10.1, tout objet
de cette catégorie admet donc une résolution droite par des objets injectifs.
Or, on vérifie aussitôt que si $(M_i)$ est un objet injectif de $I(C)$, alors
les $M_i$ sont des objets injectifs de $C$ (c'est un fait général pour les
catégories de diagrammes pour un schéma $\Sigma$ satisfaisant aux conditions
générales de prop. 1.9.2).

\begin{quote}
\small
\textbf{Correction éditoriale.} L'appel à la proposition 1.9.2 n'est pas
valide tel qu'imprimé : son assertion sur la famille de générateurs est fausse
en général. Pour $i$ dans le petit ensemble ordonné filtrant $I$, soit
$F_i:C\longrightarrow I(C)$ l'adjoint à gauche de l'évaluation en $i$, défini
par
\[
F_i(A)_j=
\begin{cases}
A,&i\leq j,\\
0,&\text{sinon,}
\end{cases}
\]
les morphismes de transition étant l'identité lorsque les deux valeurs sont
$A$. Ce foncteur est exact objet par objet. Si $E$ est un générateur de $C$,
alors $\bigoplus_{i\in I}F_i(E)$ est un générateur de $I(C)$. Puisque chaque
$F_i$ est un adjoint à gauche exact, son adjoint à droite, l'évaluation en $i$,
envoie les injectifs de $I(C)$ sur des injectifs de $C$. Ainsi $I(C)$ possède
le générateur requis par le théorème 1.10.1, les évaluations préservent les
injectifs et l'argument qui suit reste valable.
\end{quote}

Considérons alors une résolution droite de $(A_i)$ par des objets injectifs de
$I(C)$,

\LayerTarget{tohoku:eq:diagram-injective-resolution}
\[
0\longrightarrow(A_i)\longrightarrow(C_i^0)\longrightarrow(C_i^1)
\longrightarrow\cdots.
\]

\LayerTarget{tohoku:eq:component-injective-complex}
Soit, pour tout $i$, $C_i$ le complexe
\[
0\longrightarrow C_i^0\longrightarrow C_i^1\longrightarrow\cdots,
\]
on voit alors que le système $(C_i)$ répond à la question.

\LayerTarget{tohoku:eq:colimit-resolution}
Comme le foncteur $\varinjlim$ sur $I(C)$ est exact (proposition 1.8), on
obtient une résolution $C$ de $A$ par les
\[
C^p=\varinjlim_i C_i^p.
\]

\LayerTarget{tohoku:eq:T-acyclic-computation}
En vertu de la condition b), les $C^p$ sont $T$-acycliques, donc on a
\[
R^pT(A)=H^p(T(C)).
\]

\LayerTarget{tohoku:eq:T-colimit-complex}
Comme $T$ permute aux limites inductives en vertu de a), on a
\[
T(C)=\varinjlim_i T(C_i),
\]
donc, en vertu de l'exactitude du foncteur $\varinjlim$ sur $I(C')$
(résultant de l'axiome AB 5) en vertu de prop. 1.8) on obtient

\LayerTarget{tohoku:eq:cohomology-colimit-comparison}
\[
H^p(T(C))=\varinjlim_i H^pT(C_i).
\]
Or $C_i$ étant le complexe

\TohokuContinuation{170}
